\section{The Reproducibility Attrition Ladder}\label{sec:ladder}

\subsection{Why a ladder rather than a score}\label{sec:why_ladder}

Reproducibility instruments in common use return a score: a paper earns points
for releasing code, for releasing data, for pinning an environment. Scores are
convenient, but they conflate failures that have nothing to do with one another.
A paper whose repository has been deleted and a paper whose repository is
complete but whose dependencies no longer resolve both lose the same point, and
the aggregate number cannot tell the two apart. The remedies differ completely:
one is a hosting problem, the other is a packaging problem.

An ordered ladder keeps that information. Each artifact is pushed as far as it
will go, and what we record is the rung at which it stopped together with the
reason it stopped. Because the rungs are cumulative, reaching rung $n$ implies
every rung below $n$ was cleared, so the statement ``$x\%$ of papers clear rung
$n$ or better'' is unambiguous and the resulting curve is monotone by
construction. Monotonicity is not decoration: it is a property the audit script
checks on every build, because a non-monotone curve is a signal that the corpus
was scored inconsistently.

\subsection{The rungs}\label{sec:rungs}

\Cref{tab:ladder} gives the scale. Rungs 0 to 2 are properties of the manuscript
and require no execution. Rungs 3 and 4 test retrieval. Rungs 5 and 6 test
whether the declared environment can be built. Rungs 7 to 9 test execution.
Rung 10 is the scientific claim.

\begin{table}[H]
\centering
\caption{The Reproducibility Attrition Ladder. Rungs are cumulative and strictly
ordered. Rungs 0 to 9 are decided mechanically; rung 10 requires reading the paper
for a headline number and is scored only on a manually curated subset.}
\label{tab:ladder}
\small
\begin{tabular}{@{}clp{8.2cm}@{}}
\toprule
\# & Rung & Cleared when \\
\midrule
0  & \texttt{NO\_CLAIM}          & (floor) the paper makes no code-availability claim \\
1  & \texttt{CLAIM\_ONLY}        & the paper claims code exists but gives no pointer \\
2  & \texttt{URL\_PRESENT}       & a code URL appears in the paper or its metadata \\
3  & \texttt{URL\_RESOLVES}      & the URL still answers: no link rot \\
4  & \texttt{CLONED}             & the repository can be retrieved \\
5  & \texttt{DEPS\_DECLARED}     & a dependency specification exists \\
6  & \texttt{ENV\_BUILDS}        & those dependencies actually resolve and install \\
7  & \texttt{ENTRYPOINT\_FOUND}  & a runnable entry point is identifiable \\
8  & \texttt{RUNS}               & the entry point completes with exit code 0 in budget \\
9  & \texttt{PRODUCES\_OUTPUT}   & execution emits artifacts or numeric output \\
10 & \texttt{MATCHES\_PUBLISHED} & an emitted number matches a headline number \\
\bottomrule
\end{tabular}
\end{table}

Two design choices deserve defending.

First, \emph{rung 10 is never asserted by the harness}. Deciding whether an
emitted number corresponds to a paper's headline result requires reading the
paper, locating the claim, and judging whether the quantities are comparable at
all. That is human work. The harness therefore stops at rung 9, the audit script
verifies that no record claims rung 10 automatically, and the rung-10 figure is
reported only for a manually curated subset and never extrapolated to the
corpus.

Second, \emph{network access is denied during execution and permitted during
installation}. A repository that downloads its data when it runs is not
self-contained: a reader without the author's credentials, or reading after the
endpoint moves, cannot reproduce the result. Allowing the fetch would let our own
connectivity paper over that. Installation is a different matter, since package
indexes are the normal mechanism by which environments are built, so the
network is available there.

\subsection{A closed vocabulary for failures}\label{sec:reasons}

Recording \emph{that} an artifact failed at rung 6 is less useful than recording
\emph{why}. Free-text error capture makes that unaggregatable, so each rung
carries a fixed set of admissible reasons and the data structure refuses any
value outside it. The result is that the failure mix at each rung is a table.

The vocabulary separates causes that call for different responses. At rung 6, for
instance, \texttt{package\_not\_found} (a dependency has left the index),
\texttt{python\_\allowbreak version\_\allowbreak unsupported} (the pin has aged
out of the supported range), and
\texttt{build\_\allowbreak from\_\allowbreak source\_\allowbreak failed} (a wheel is unavailable for the
platform) are distinct, and only the first is really about the paper's own
choices.

One category is reserved for our own defects. If the harness locates a
dependency specification and then fails to read it, that is a bug in the
measurement, not a property of the paper, and it is recorded as
\texttt{spec\_unreadable\_by\_backend} and excluded from every claim about the
literature. This is not hypothetical: an early version of the container backend
passed a bare filename for specifications found in a subdirectory, and the
resulting installer error was initially classified as a dependency conflict, that
is, as evidence against the papers. The separate code, and a regression test
pinning it, exist because that failure mode is easy to reintroduce and
impossible to notice in an aggregate.

\subsection{Execution environment}\label{sec:environment}

Each repository is cloned to a scratch directory, its declared environment is
built, and its entry point is executed inside a container with a fixed CPU and
memory allocation. Containers matter for two reasons: dependency isolation, so
that one repository's pins cannot affect another's, and process isolation, since
we are running code written by authors who never anticipated a stranger
executing it. Where a container runtime is unavailable the harness falls back to
a per-repository virtual environment combined with kernel-level sandboxing, which
confines filesystem writes and denies network but shares the host kernel. The
backend used is recorded on every result, and results produced under different
isolation regimes are reported separately rather than pooled.

Entry-point selection is the most discretionary step in the pipeline, so it is
run as a switch rather than a constant. The \emph{convention} route looks for
conventionally named scripts and otherwise accepts a single unambiguous top-level
script; build scripts and configuration modules are excluded, because executing
\texttt{setup.py} exercises the packaging toolchain rather than the paper's
analysis. The \emph{documentation} route additionally reads the repository's
README and takes the command the authors themselves give, subject to three
restrictions: the command must name a target that exists in the repository, it
must not be an installation or testing step, and it must not be a compound shell
pipeline, since those cannot be reproducibly parsed.

Both routes exclude directories whose contents are never the analysis: test
suites, example folders, documentation, and vendored dependencies. That exclusion
is not hypothetical tidiness. Two repositories in this corpus ship a copy of a
Chinese-tokenization library, and its bundled \texttt{test/demo.py} satisfied the
naming heuristic; the harness executed the library's demo and recorded its output
as though the paper's pipeline had produced it. Both were counted at the top of
the ladder until the outputs were read and found to be identical to each other
and unrelated to either paper. Running a repository's tests is also not running
its analysis, which is why a documented \texttt{pytest} invocation is refused as
an entry point.

The exclusion is weaker than that principle implies, and the gap is worth stating
rather than leaving for a reader to find. It operates on directory names and on
the literal \texttt{pytest} command, not on file names, so a test script sitting
at the repository root remains selectable. Nine of the 71 entry points found
under the primary specification are test artifacts of this kind, including two of
the seven repositories that reach the top of the ladder. We report the
consequence in \Cref{sec:rung10_results} rather than tightening the rule
retroactively: it is the clearest evidence in the study that clearing rung 9 is a
weaker achievement than its name suggests, and suppressing it would have made the
harness look better and the finding weaker.

Both configurations are run over the whole corpus and both curves are reported
(\Cref{sec:sensitivity}). Neither matches a determined human who would read the
source and infer the intended invocation, so both are lower bounds; reporting the
pair makes the size of that discretion visible instead of burying it in a design
choice. Where no entry point is identified under the active configuration, the
harness records \texttt{no\_entrypoint\_identifiable} rather than guessing: a
harness that invents plausible commands would report their failure as the
paper's.

\subsection{What the ladder does not measure}\label{sec:ladder_limits}

An artifact that stops low on the ladder is not thereby wrong. It is not
independently executable as released, on a current platform, within our time
budget, by a reader who did not write it. That is a narrower claim than
``irreproducible'' and it is the only claim the instrument supports. Conversely,
clearing rung 9 shows that code ran and produced output; it says nothing about
whether the output is correct, whether the method is sound, or whether the
finding replicates on new data. Rung 10 speaks to the first of those and to
nothing beyond it.
